% SPDX-FileCopyrightText: Copyright (c) 2026 NVIDIA CORPORATION & AFFILIATES. All rights reserved.
% SPDX-License-Identifier: Apache-2.0
%
% ROLE: the tuning details that sec:training summarizes: the harness's objective bookkeeping, every
%   search space (app:spaces; merges training.tex 79-265 with baselines.tex 95-233 of the base
%   revision <commit-50>, so tab:baseline-spaces absorbs the former tab:spaces), the rank-stability
%   diagnostic, alg:train, the initial-point fits, the other learned arms' spaces, the
%   sign-constraint clamp, the drift check (sec:drift, merging training 543-562 and model 496-521)
%   and compute and numerics (sec:compute). Facts unchanged by the move.
% EVIDENCE: WT/benchmarks/learned_routing/learned_routing/{train.py,space.py,train_cli.py};
%   WT/benchmarks/learned_routing/remote/{p2orch.py,node_train.sh}; CR/runs/phase2/spaces/*.yaml;
%   CR/runs/phase2/MANIFEST.json; CR/runs/phase2/scripts/{flipcal.py,bc_fit.py};
%   CR/facts/gate.json full_plan; CR/facts/pilot.json; CR/facts/phase2_prep.json;
%   CR/facts/tierA.json; CR/facts/tierBC.json; CR/facts/remote.json; CR/audits/phase2-mid/fix.md;
%   CR/CONTRACT.md A9-A12, A15-A17; CR/literature/LESSONS.md LR-03, LR-05, LR-07, LR-15;
%   cma 4.5.0 in WT/.venv (cma/options_parameters.py).
% NAMING: "lr-train" is "the tuning harness" (prov: lr-train); "phase 2" is "the main tuning";
%   tiers A/B/C are described by their content (prov comments mark each place).

\section{Tuning details}
\label{app:tuning}

This appendix holds the tuning details that \cref{sec:training} summarizes: how the tuning harness
keeps scores paired, every search space and how its step size was set, the diagnostic that is
logged but steers nothing, one tuning run in pseudocode, the initial points and the other learned
arms, how the sign constraint is implemented, the drift check that preceded M2, and the compute
behind the runs.

\subsection{Objective bookkeeping}

% src: train.py (objective_value, pair_score, eval_set)
Before scoring, the tuning harness checks that every returned record carries its task's cell,
replicate and policy hash, and that candidate and reference cover the same (cell, replicate) set;
a mismatch stops the run rather than producing an unpaired score. After its first replay,
$\piref$'s record comes from the result cache.
% prov: lr-train

% src: train.py Trainer.run ("worst = min(finite) - 1.0"), Trainer.score (reference error raises);
%      CR/facts/gate.json criteria.1_error_rate.lr_train_candidate_failures = 0 (4 runs x 480 evals)
A candidate whose replay fails on any cell receives the generation's lowest finite objective minus
one, so CMA-ES ranks it last; a failed reference replay stops the generation, which is retried on
the next call because errors are not cached. In the pilot no candidate failed in 4 runs of 480
evaluations (\emph{preliminary (pilot, train subset)}).
% src: 34 cells x 2 replicates x (16 candidates + reference) = 1,156 replays; re-evaluation
%      34 x 2 x (2 + 1) = 204 (arithmetic; matches CR/facts/phase2_prep.json early_progress
%      replays_fresh 1156 and tasks_requested 1360 at generation 1)
% src: CR/facts/pilot.json measurements.seconds_per_evaluation.full_train_local
%      replay_s_per_candidate_evaluation 700.895
One generation replays 16 candidates and $\piref$ on 34 cells $\times$ 2 replicates, \num{1156}
replays, plus 204 for the rank-stability diagnostic (\cref{app:rank-diagnostic}); a single
candidate evaluation on the full train split costs about 700.9 replay-seconds, as measured in the
pilot.

\subsection{Search spaces}
\label{app:spaces}

% src: WT/benchmarks/learned_routing/learned_routing/space.py (module docstring, Coord, cma_bounds);
%      every CR/runs/phase2/spaces/*.yaml param has bounds (checked 2026-10-03)
\paragraph{Internal coordinates.}
Each policy's \code{space.yaml} lists its free parameters with bounds and a linear or log scale,
its pinned parameters, and CMA-ES options. CMA-ES searches internal coordinates $\searchvec$: each
bounded parameter maps to $[0,1]$, linearly or in log space, and integer parameters are rounded
after decoding. pycma's boundary transform keeps every sample inside the box, so every candidate
decodes to a valid policy file. Every parameter of the main tuning is bounded.
% prov: "Every phase-2 parameter is bounded."

\paragraph{M1's bounds.}
% src: CR/facts/pilot.json feature_scaling.pooled_within_set_sd (geomean over the 12 pilot-subset
%      cells of each cell's mean within-set sd; CR/runs/pilot/screen/feature_scales.json "pooled");
%      ratios = sd(f0)/sd(f): 25.6, 17.7, 5.33, 80.8, 109.6, 21.2, 3.22 (also flip_rate.json s_j);
%      bounds from CR/runs/phase2/spaces/m1_learned_v1.yaml
A feature's scale ratio $\mathrm{sd}(x_{\cdot,0}) / \mathrm{sd}(x_{\cdot,f})$ is the coefficient
at which feature $f$ varies across a candidate set as much as the anchor does; at its bound of
about four times that ratio, a feature's term spreads four times as widely as the anchor's
(\cref{tab:m1-space}). The standard deviations are measured within candidate sets, as the geometric
mean over the pilot's 12 train cells of each cell's mean, on candidate tables reconstructed from
$\piref$'s per-request rows, because replay does not log candidate tables. The reconstruction is
approximate (router load is rebuilt from arrivals, first tokens and completions, and overlap is
exact only for the chosen worker and for earlier turns of the same session), so the ratios set
scales, not policy inputs. The initial placeholder bounds of $\pm 8$ could not have let
\code{kv\_load\_frac} or \code{active\_requests\_s} matter.
% prov: "The build stage's placeholder bounds of +-8"
% src: space.py (bounded linear: u = (x - low)/(high - low)); sigma0 0.004 -> physical std
%      0.004 x 2 x bound = 0.032 x ratio (m1_learned_v1.yaml comment); LESSONS.md LR-05
%      (Hansen tutorial App. A, pdf p.29; ARS sec. 3.2, pdf p.7-8)
Because each bound is proportional to the feature's scale ratio, one internal unit moves every
coefficient by the same multiple of its ratio, so the initial search distribution is isotropic in
units of decision-relevant spread. This is the per-variable scaling \citet{hansen2016tutorial}
recommends when initial ranges differ by orders of magnitude, and the analogue of the state
normalization without which \citet{mania2018ars} could not train a linear policy on their
hardest task.

\begin{table}[tbp]
  \centering\small
  \caption{M1's search space and its three initial points. Within-set sd: geometric mean over the
    pilot's 12 train cells of the mean within-candidate-set standard deviation, on reconstructed
    tables. Ratio: $\mathrm{sd}(x_{\cdot,0})/\mathrm{sd}(x_{\cdot,f})$. The search box is
    $|\theta_f| \le$ bound. Start s1 is $\theta = -\basis{0}$ (all free coefficients 0); s2 and s3
    are described in \cref{sec:restarts,app:starts}.}
  \label{tab:m1-space}
  % src: CR/facts/pilot.json feature_scaling.pooled_within_set_sd; CR/runs/pilot/screen/flip_rate.json s_j;
  %      CR/runs/phase2/spaces/m1_learned_v1.yaml bounds; CR/facts/phase2_prep.json m1_inits.s2.theta,
  %      m1_inits.s3.theta (s2 rounded to 3 significant digits)
  \begin{tabular}{@{}rlrrrrr@{}}
    \toprule
    $f$ & Feature & Within-set sd & Ratio & Bound & s2 & s3 \\
    \midrule
    0 & \code{default\_logit\_scaled}     & 6.607   & 1     & pinned & $-1$     & $-1$ \\
    1 & \code{overlap\_frac}              & 0.2585  & 25.6  & 100    & 63.9     & 76.8 \\
    2 & \code{new\_prefill\_tokens\_k}    & 0.3742  & 17.7  & 70     & 4.18     & 0 \\
    3 & \code{active\_prefill\_tokens\_k} & 1.239   & 5.33  & 21     & 1.28     & 0 \\
    4 & \code{kv\_load\_frac}             & 0.08173 & 80.8  & 320    & 12.5     & 0 \\
    5 & \code{active\_requests\_s}        & 0.06031 & 110   & 440    & $-138$   & 0 \\
    6 & \code{session\_affinity}          & 0.3111  & 21.2  & 85     & $-0.324$ & 42.4 \\
    7 & \code{isl\_x\_prefill\_load}      & 2.053   & 3.22  & 13     & $-0.196$ & 0 \\
    \bottomrule
  \end{tabular}
\end{table}

\subsubsection{Baseline spaces}
\label{sec:baseline-spaces}

% src: CR/runs/phase2/spaces/*.yaml (ranges, scales, inits, pins, sigma0, maxstd);
%      CR/facts/phase2_prep.json decisions[1..3], spaces.*.free/fixed/base_parameters
Each baseline searches the numeric parameters its catalog entry exposes, over ranges set in the
parameter's own units, log-scaled where a parameter spans orders of magnitude, starting from its
shipped defaults (\cref{tab:baseline-spaces}); restarts 1, 2 and 3 use CMA-ES seeds 1, 2 and 3.
Four kinds of parameter are pinned rather than searched.
\begin{itemize}
  \item \textbf{Rescalings.} A parameter that only rescales others adds a direction along which
    decisions do not change, and CMA-ES's step size drifts along such directions
    \citep{hansen2016tutorial}. Three are pinned. \policy{ramjet}'s affinity block size stays at
    512 tokens: multiplying \code{affinity\_block\_tokens} by some factor while dividing
    \code{max\_affinity\_blocks} and \code{alpha} by the same factor leaves its decisions unchanged
    up to integer rounding. \policy{llm-d-precise-prefix}'s prefix weight stays at 2: scaling all
    three of its weights never changes its choice. And \policy{llm-d-optimized-baseline}'s peak
    prefill rate stays at \num{15928} tokens/s: with the throughput source, its gate compares a
    token difference divided by \code{peak\_prefill\_tokens\_per\_second} with
    \code{max\_ttft\_penalty\_ms}, so their product is the only quantity that policy identifies.
  \item \textbf{Categorical choices.} CMA-ES searches continuous coordinates, so these structural
    parameters keep their shipped values: \policy{lmetric}'s hot-spot detector stays on,
    \policy{ramjet}'s affinity basis stays \code{relative} and the optimized baseline's TTFT source
    stays \code{throughput}.
  \item \textbf{Capacity.} The sticky-session map keeps its \num{65536}-session bound.
  \item \textbf{Inputs replay does not supply.} \policy{llm-d-precise-prefix}'s per-class weights
    stay empty, because replay sets no policy class.
\end{itemize}
% src: CR/facts/phase2_prep.json decisions ("only sticky-session reads the KvRouterConfig default-cost
%      knobs ... verified: sticky@x0 == M0@x0 and sticky@hi == M0@hi on a session-free Mooncake cell;
%      no other baseline reads a tunable KvRouterConfig knob; router_queue_threshold stays for ablation a")
Of the baselines, only \policy{sticky-session} reads the default cost's router-level knobs, because
its first-turn and fallback choice is the default cost function; both sticky variants therefore
also search the five M0 knobs over M0's ranges. Sticky-session and M0 at the same knob values gave
identical results on a session-free Mooncake cell, both at the start point and at the upper corner
of the box. No other baseline reads a tunable router-configuration knob, and
\code{router\_queue\_threshold} keeps its default for every policy; the queue-threshold ablation
tunes it separately (\cref{sec:ablations}).

% src: CR/runs/phase2/spaces/m0_default_tuned.yaml (comment, bounds, init); CR/facts/gate.json
%      full_plan.search_spaces.M0; CR/facts/pilot.json measurements.cma_progress_curves ("prefill_load_scale
%      moved to 17-25 of a 50 upper bound")
M0's space is the pilot's, with the prefill-load scale $\bpf$ widened to $[0.05, 500]$ on a log
scale because the pilot's tuned values (17 to 25) sat near the top of the old range $[0.05, 50]$.
The decode request weight $\breq$ starts at 0.5, the bottom of its log range $[0.5, 4096]$, which
stands in for the shipped value 0.

% src: CR/facts/finalists.json headline.val_best_baseline (ramjet@p2-ramjet-s3-g25-best_so_far);
%      CR/report/REPORT.md section 3.1 "Selected config (val k0-2)" column
\Cref{tab:leaderboard-full} gives each tuned baseline's selected restart and checkpoint with its
validation and test scores; the selected \policy{ramjet}, for instance, came from restart 3 at
generation 25, with \code{alpha} 0.883, \code{max\_affinity\_blocks} 68 and
\code{load\_unit\_tokens} 4552 (\cref{fig:api}).

\begin{table}[tbp]
  \centering\footnotesize
  \caption{Tuning spaces. Each searched parameter is listed with its range, ``log'' if it is
    searched on a log scale, and its starting value, which is the shipped default except where
    marked. $d$: free parameters. Flips: the share of decisions that one CMA-ES step of size
    $\sigmazero$ changes, measured from the policy's start on candidate tables reconstructed
    from the pilot's default-router replays of 12 train cells (target 0.175). Every baseline space
    caps the per-coordinate standard deviation at 0.3. M0's symbols are the router knobs
    \code{overlap\_score\_credit} ($\bov$), \code{overlap\_score\_credit\_decay} ($\bdecay$),
    \code{prefill\_load\_scale} ($\bpf$) and \code{decode\_active\_request\_weight} ($\breq$).
    M1's row is for comparison; M1 has three starts (\cref{tab:m1-space}).}
  \label{tab:baseline-spaces}
  % src: CR/runs/phase2/spaces/{m0_default_tuned,two_tier,lmetric,ramjet,dualmap,chwbl,
  %      llmd_precise_prefix,llmd_optimized_baseline,sticky_hard,sticky_bounded,m1_learned_v1}.yaml;
  %      CR/facts/phase2_prep.json spaces.<file>.{dimension,sigma0,free,fixed,base_parameters,
  %      flip_calibration.flip_at_sigma0_interpolated}
  %      (0.16938, 0.08297, 0.00058, 0.13332, 0.01828, 0.17260, 0.17698, 0.03404, 0.17908, 0.17460),
  %      rounded to 3 places (lmetric to 1 significant digit); M1 flips: CR/facts/pilot.json decisions
  %      ("sigma0 0.004 internal ... about 20% decision flips vs theta0"),
  %      feature_scaling.m1_flip_frac_by_c (measured against theta = -e0)
  % Merge note: the former tab:spaces (training.tex 211-246 at the base revision) had the same rows,
  %   dimensions, sigma0 and flips; its code names for the llm-d-precise-prefix weights are used below.
  %   The caption's "Every baseline space" was "Every space" (M1's spaces set no explicit cap;
  %   CR/runs/phase2/spaces/*.yaml: maxstd 0.3 in every baseline space, absent in the three M1 spaces).
  \begin{tabularx}{\linewidth}{@{}P{0.17\linewidth}YP{0.2\linewidth}ccc@{}}
    \toprule
    Policy & Searched: range (scale), start & Pinned & $d$ & $\sigmazero$ & Flips \\
    \midrule
    default cost (M0)
      & $\bov$ $[0, 4]$, 1;\newline $\bdecay$ $[0, 4]$, 0;\newline $\bpf$ $[0.05, 500]$ log, 1;\newline
        $\breq$ $[0.5, 4096]$ log, 0.5\textsuperscript{a};\newline \code{router\_temperature} $[0, 1]$, 0
      & --- & 5 & 0.11 & 0.169 \\
    \addlinespace
    two-tier
      & \code{cache\_threshold} $[0, 1]$, 0.5;\newline \code{balance\_abs\_threshold} $[1, 1024]$ log, 32;\newline
        \code{balance\_rel\_threshold} $[1, 8]$ log, 1.1
      & --- & 3 & 0.3 & 0.083 \\
    \addlinespace
    LMetric
      & \code{class\_prefix\_blocks} $[1, 64]$ log, 4;\newline \code{window\_requests} $[50, \num{50000}]$ log, 1000
      & hot-spot detector on & 2 & 0.3 & 0.0006 \\
    \addlinespace
    Ramjet
      & \code{alpha} $[0.1, 160]$ log, 4;\newline \code{max\_affinity\_blocks} $[1, 1024]$ log, 32;\newline
        \code{load\_unit\_tokens} $[256, \num{262144}]$ log, 8192
      & affinity block 512 tokens\textsuperscript{b};\newline basis \code{relative} & 3 & 0.3 & 0.133 \\
    \addlinespace
    DualMap
      & \code{hash\_prefix\_blocks} $[1, 64]$ log, 4\textsuperscript{c};\newline
        \code{pending\_prefill\_token\_budget} $[2048, \num{2097152}]$ log, \num{65536};\newline
        \code{window\_requests} $[50, \num{50000}]$ log, 1000
      & --- & 3 & 0.3 & 0.018\textsuperscript{c} \\
    \addlinespace
    CHWBL
      & \code{prefix\_tokens} $[16, \num{16384}]$ log, 256\textsuperscript{c};\newline \code{load\_factor} $[1, 8]$ log, 1.25
      & --- & 2 & 0.16 & 0.173\textsuperscript{c} \\
    \addlinespace
    llm-d precise prefix
      & \code{weights.queue} $[0.01, 100]$ log, 1;\newline \code{weights.kv\_cache\_utilization} $[0.01, 100]$ log, 1;\newline
        \code{prefix\_match\_length\_weight} $[0, 1]$, 0;\newline
        \code{prefix\_match\_length\_scale\_tokens} $[512, \num{131072}]$ log, 8192
      & \code{weights.prefix} 2\textsuperscript{b};\newline per-class weights empty & 4 & 0.16 & 0.177 \\
    \addlinespace
    llm-d optimized baseline
      & \code{affinity\_threshold} $[0, 1]$, 0.8;\newline \code{max\_ttft\_penalty\_ms} $[100, \num{3240000}]$ log, \num{18000};\newline
        \code{queue\_threshold\_tokens} $[\num{16384}, \num{1073741824}]$ log, \num{4194304}
      & peak prefill rate \num{15928} tokens/s\textsuperscript{b};\newline TTFT source \code{throughput}
      & 3 & 0.3 & 0.034 \\
    \addlinespace
    sticky, hard
      & the five default-cost knobs, with M0's ranges and starts
      & mode \code{hard};\newline \num{65536} sessions & 5 & 0.17 & 0.179 \\
    \addlinespace
    sticky, bounded
      & \code{load\_factor} $[1, 8]$ log, 1.25;\newline the five default-cost knobs, as M0
      & mode \code{bounded};\newline \num{65536} sessions & 6 & 0.067 & 0.175 \\
    \midrule
    learned (M1)
      & $\theta_1,\dots,\theta_7$, bounds in \cref{tab:m1-space}
      & $\theta_0 = -1$, $\temp = 0$ & 7 & 0.004 & about 0.20\textsuperscript{d} \\
    \bottomrule
  \end{tabularx}

  \smallskip
  \begin{minipage}{\linewidth}\footnotesize
    \textsuperscript{a}\,The shipped value of $\breq$ is 0, which a log scale cannot reach; 0.5, at
    most half a block per active request, stands in for it.
    \textsuperscript{b}\,Pinned because it only rescales the other parameters.
    \textsuperscript{c}\,Hashed key depth, held at its default when $\sigmazero$ was set: any
    change of it redraws the prefix-to-worker map (0.69 of decisions change for \policy{chwbl} and
    0.18 for \policy{dualmap} at a step of 0.03).
    \textsuperscript{d}\,Measured in the pilot, from $\theta = -\basis{0}$.
  \end{minipage}
\end{table}

\subsubsection{Initial step size}

% src: CR/runs/phase2/scripts/flipcal.py (docstring: SIGMAS, SAMPLES 8, TARGET 0.175, TV distance,
%      stateful policies, clipping); CR/facts/phase2_prep.json flip_calibration (rule, cells = the 12
%      pilot-subset cells, target 0.175); CR/facts/DEVIATIONS.md 2026-10-03 phase-2 prep "sigma0 for tuned baselines"
Parameter units are not comparable across policies, so $\sigmazero$ is set in decision space.
For a step size $\sigma$, we draw internal points from $\mathcal{N}(\searchvec_0, \sigma^2 I)$
clipped to the box, decode them, and measure the share of routing decisions that differ from the
policy's own decisions at its start $\searchvec_0$, as the total-variation distance between the
two decision distributions. The decisions are one-step counterfactuals on the candidate tables
reconstructed from the pilot's default-router replays of its 12 train cells, with 8 draws per
$\sigma$ and cell; stateful policies run their own state machine over each cell in arrival order.
$\sigmazero$ is the $\sigma$, interpolated on a grid from 0.03 to 0.35, at which this flip rate
reaches 17.5\%, so every policy starts with a comparable amount of exploration measured in
decisions. The literature review's suggested band of 10--30\% is itself a hypothesis.
% prov: LR-05
% src: CR/facts/pilot.json feature_scaling.m1_flip_frac_by_c (c = 0.03 -> 0.196, c = 0.05 -> 0.239;
%      sigma0 0.004 corresponds to c = 0.032); CR/runs/phase2/spaces/m1_learned_v1.yaml comment ("~20%")
For M1, $\sigmazero = 0.004$ internal units, a standard deviation of 0.032 scale ratios per
coefficient, flips about 20\% of decisions relative to $\theta = -\basis{0}$.
% src: CR/facts/phase2_prep.json spaces.<name>.flip_calibration.flip_at_sigma0_interpolated,
%      target_reached_by_0.3; flip_calibration.rule (key depth: "flips 0.69 / 0.18 already at sigma 0.03");
%      CR/facts/STATE.md phase2-prep ("unreachable by 0.3 -> 0.3: two-tier 0.083, lmetric 0.0006,
%      ramjet 0.133, dualmap 0.018, optimized-baseline 0.034")
Five baseline spaces cannot reach the target by $\sigma = 0.3$, the step-size cap below, and start
at 0.3: two-tier (0.083 of decisions change), \policy{lmetric} (0.0006), \policy{ramjet} (0.133),
\policy{dualmap} (0.018) and the optimized baseline (0.034). Around their shipped defaults their
decisions barely respond to their parameters: \policy{lmetric}'s flip rate at $\sigma = 0.3$ is
0.06\%. For \policy{chwbl}'s and \policy{dualmap}'s hashed prefix depths, any change re-draws the
whole prefix-to-worker map, so the objective is flat between block boundaries and jumps across
them (flip rates of 0.69 and 0.18 already at $\sigma = 0.03$); those two coordinates are held at
their defaults during the measurement and share their space's $\sigmazero$.

\subsubsection{Step-size cap}

% src: CR/runs/phase2/spaces/*.yaml (cma maxstd 0.3 in every baseline space; absent in the three M1
%      spaces); CR/facts/pilot.json measurements.cma_progress_curves ("seed 1's sigma then diverged
%      (0.18 -> 10.8 internal units ...): flat directions in decode_active_request_weight/temperature"),
%      recommendations_for_phase2[1]
% src: cma 4.5.0 cma/options_parameters.py:84-85 (maxstd None, maxstd_boundrange 1/3) and
%      cma/evolution_strategy.py:1179-1183, 1539-1564 (_stds_into_limits rescales sigma_vec, not sigma);
%      checked 2026-10-03: CMAEvolutionStrategy with bounds [0,1]^7 has opts['maxstd'] = 1/3 per coordinate
% check: whether the pilot's M0 seed-1 per-coordinate stds actually reached the 1/3 cap (history logs only sigma)
Every baseline space caps pycma's per-coordinate standard deviation (\code{maxstd}) at 0.3 internal
units. The pilot recommended the cap after M0's first run, whose global step size grew from 0.18 to
10.8 internal units after its validation score had plateaued, along nearly flat directions in the
decode request weight and the temperature (\emph{preliminary (pilot, train subset)}). M1's space,
whose step size stayed small in the pilot, sets no explicit cap. Two details qualify the measure.
pycma 4.5.0 already caps each bounded coordinate at one third of its range by default, so on these
$[0,1]$ coordinates the explicit cap tightens the default from $1/3$ to 0.3, and the default
applies to M1. The cap also limits the per-coordinate spread of the samples, not the global step
size $\sigma$, which can still drift.

\subsection{Rank-stability diagnostic}
\label{app:rank-diagnostic}

% src: train.py Trainer._reeval (count = max(2, ceil(reeval_frac x popsize)); fresh_ks =
%      (g K + K + j) mod P; "diagnostic"); es.tell uses only the K = 2 training values;
%      CR/runs/phase2/MANIFEST.json args --reeval-frac 0.125; p2orch.py budget_report comment;
%      CR/literature/LESSONS.md LR-03 (UH-CMA-ES sec. IV, pdf p.8-9: r = max(0.1, 2/lambda));
%      CR/facts/pilot.json runs.{m0-s1,m0-s2,m1-s1,m1-s2}.uh_cmaes_top2_swaps (8, 11, 7, 11 of 30)
After each update, the two best candidates are re-scored on the same cells with the next
generation's replicates $\Kset{\gen+1}$, and the change in their order is logged. The re-scored
share, $2/\popsize = 0.125$, equals UH-CMA-ES's default rate $\max(0.1, 2/\popsize)$
\citep{hansen2009uh}, but unlike UH-CMA-ES nothing reacts to the measurement: neither the step size
nor the number of evaluations per candidate changes. It runs identically for every policy, does not
count against $\budget$, and is skipped in a generation whose wall-clock chunk ends during it. In
the pilot, the top two candidates swapped in 7 to 11 of 30 generations per run
(\emph{preliminary (pilot, train subset)}).

\subsection{One tuning run}

\Cref{alg:train} summarizes one restart of the tuning harness.

\begin{algorithm}[tbp]
\caption{Tuning one policy (one restart of the tuning harness). Every policy, learned or baseline,
  runs this with the same $\budget$, $\popsize$, $\Krep$, cells, validation schedule and selection
  rule.}
\label{alg:train}
% prov: "Tuning one policy (\code{lr-train}, one restart)"
\begin{algorithmic}[1]
\Require space (bounds, pins, $\sigmazero$, cap), start $\searchvec_0$, seed, cells $\Cset{train}$
  and $\Cset{val}$
\State initialize CMA-ES at $\searchvec_0$ with step size $\sigmazero$ and population $\popsize$
\State $\textit{best} \gets \bot$; \quad $\textit{selected} \gets \bot$
\For{$\gen = 0, 1, \dots, \budget/\popsize - 1$}
  \State $\Kset{\gen} \gets \{(\gen\Krep + i) \bmod \Kpool : i < \Krep\}$
  \State sample candidates $\searchvec_1,\dots,\searchvec_{\popsize}$; checkpoint \Comment{\code{ask}}
  \State replay every candidate and $\piref$ on $\Cset{train} \times \Kset{\gen}$; check CRN coverage
  \State compute $\Jobj$ per candidate (\cref{eq:objective}); a failed one gets the worst finite $\Jobj$ $-1$
  \State update CMA-ES with fitness $-\Jobj$ \Comment{\code{tell}}
  \State $\textit{best} \gets$ the candidate with the highest $\Jobj$ seen so far
  \State re-score the top 2 on $\Cset{train} \times \Kset{\gen+1}$; log rank changes \Comment{diagnostic only}
  \If{$(\gen + 1) \bmod 5 = 0$}
    \State score \textit{best} and the CMA-ES mean on $\Cset{val} \times \{0,1,2\}$
    \State $\textit{selected} \gets$ the validated candidate with the highest score so far
  \EndIf
  \State checkpoint (CMA-ES state and random state)
\EndFor
\State \Return \textit{selected}, and \textit{best} for reference
\end{algorithmic}
\end{algorithm}

\subsection{Restarts and initial points}
\label{app:starts}

% src: CR/facts/phase2_prep.json decisions[0], m1_inits.restart_to_wave; CR/runs/phase2/MANIFEST.json
%      (seed = restart, wave = restart for tier A); CR/facts/gate.json full_plan.scheduling
Restart 1 of every policy runs in the first wave, restart 2 in the second and restart 3 in the
third, so a compute cut at any point removes the same restart from every policy.

\paragraph{The behavior clone (s2).}
% src: CR/runs/phase2/scripts/bc_fit.py (docstring: model, pinning, bounds, L-BFGS-B, equal cell
%      weights, stride sampling); CR/facts/phase2_prep.json m1_inits.s2_fit (decisions 144048, cells 34,
%      agreement_in_sample 0.8748, agreement_cv_2fold_by_cell 0.8659/0.8827, agreement_theta0_default_rule
%      0.5121, at_bound []); CR/facts/pilot.json mde.learned_vs_val_best_default_arm_heuristic.comparator
Start s2 clones \policy{llm-d-precise-prefix}, the best of the untuned heuristics evaluated on
validation in the pilot. It is the conditional-logit maximum-likelihood fit to that heuristic's
choices: with $\Pr[i \mid \cand] \propto \exp(\theta^{\top}\fvec_i)$ over reconstructed candidate
tables, it maximizes the mean over cells of the mean log-probability of the worker the heuristic
chose, with $\theta_0$ pinned at $-1$ and $\theta_1,\dots,\theta_7$ held inside M1's box. The
problem is convex and bounded and is solved by L-BFGS-B. The data are the heuristic's own replays
at shipped defaults on all 34 train cells, replicate 0: \num{144048} decisions in a stride sample
of each cell. The fit chooses the heuristic's worker on 0.875 of these decisions (0.866 and 0.883
in a two-fold split by cell), against 0.512 for the default rule $\theta = -\basis{0}$, and no
coefficient sits at its bound.

\paragraph{The cache-heavy point (s3).}
% src: CR/facts/phase2_prep.json m1_inits.s3 (theta, why), decisions[7]; CR/facts/pilot.json
%      validation_vs_best_heuristic.llm-d-precise-prefix.M1_tuned.by_family_mean_clipped_log_ratio.agentx
%      (-0.01898); CR/literature/LESSONS.md LR-07 action 3
Start s3 sets $\theta_1 = 3 \times 25.6 = 76.8$ and $\theta_6 = 2 \times 21.2 = 42.4$, so that
across a candidate set the overlap term spreads three times and the session-affinity term twice as
widely as the anchor, with all other free coefficients 0, so load enters only through the anchor.
It targets the agent workloads on which the pilot's M1 trailed the cache-aware heuristic (by
$-0.019$ in mean clipped log-ratio on AgentX, \prelim). It was preferred to a fit to the
second-best heuristic family because \policy{lmetric}'s product of uncached tokens and load is not
representable with \code{v1} features, and a fit to \policy{ramjet} would be another linear
cache-minus-load score close to s2.

\paragraph{How the starts differ.}
% src: CR/facts/phase2_prep.json m1_inits.{s2,s3}.one_step_flips_vs_theta0 (0.4645, 0.4198),
%      s2_vs_s3_one_step_disagreement (0.1096), s3.agreement_with_precise_prefix (0.8722),
%      s2_fit.caveat, s2_fit.sensitivity_hash_aware_overlap (agreement 0.7847, theta1 40.3),
%      s2_agreement_by_family (agentx 0.797, mooncake 0.914, sessions 0.865)
On one-step decisions, s2 and s3 differ from the default rule on 46\% and 42\% of requests and from
each other on 11\%. The clone's agreement should be read with the reconstruction's bias in mind:
overlap is exact only for the chosen worker, which favors overlap-heavy rules. With a hash-aware
reconstruction the fit agrees on 0.785 of decisions and its overlap coefficient is 40 instead of
64, and the hand-set point s3 agrees with the heuristic on 0.872, almost as often as the fit.
Agreement is lowest on AgentX (0.797, against 0.914 on Mooncake and 0.865 on sessions).

\subsection{Other learned arms}

% src: CR/runs/phase2/MANIFEST.json (m2r2-*, m1cont-* tier B: warm start (theta*_M1, p = 0), context
%      CMA_stds about 0.3x theta's; budget 400, 3 seeds); CR/CONTRACT.md A15 (M1-noaff: theta_6 pinned 0,
%      same multi-start inits, 3 restarts at B = 400, post-hoc theta_6 = 0 check on val k0-2), A17.1 (M1-v2 inits)
Each arm of \cref{tab:ladder} is trained like M1, with $\budget = 400$, three restarts, the same
cells and the same selection rule. Their spaces and starting points are as follows.
\begin{itemize}
  \item \textbf{M2 rank 2 and M1 continued.} M2 starts at the selected M1 coefficients with $p = 0$
    and context step sizes 0.3 times $\theta$'s; its two named sources, \code{isl\_k} and
    \code{mean\_active\_prefill\_tokens\_k}, were fixed before any M2 evaluation, and $\theta_7$ is
    pinned at 0 because the (\code{isl\_k}, feature 3) context entry duplicates it. It is compared
    with M1 continued from its selected point for the same added 400 evaluations.
    % prov: LR-15 (compare a nested model with its base continued for the same added budget)
  \item \textbf{M1-noaff.} M1's three starts with $\theta_6$ pinned at 0; a further check, without
    retraining, zeroes $\theta_6$ in the selected M1.
    % prov: Amendment A15
  \item \textbf{M1-v2.} Coefficients 0 to 7 keep M1's bounds, so M1-v2 nests M1; coefficients 8 to
    22 are bounded at four times their scale ratio, measured on candidate tables reconstructed from
    the default router's train replays. $\sigmazero = 0.0028$ flips 19.6\% of decisions at the s1
    start, the rate M1's space has at its own $\sigmazero$. The starts are s1, $\theta = -\basis{0}$;
    s2, the behavior clone of \policy{llm-d-precise-prefix} refit on \code{v2} features (with
    \code{hash\_home} pinned at 0 in the fit, because its offline reconstruction is approximate);
    and s3, \policy{lmetric}'s representable point
    $\theta = -\basis{0} - 20.7\,(\basis{\text{\code{log\_ptok}}} + \basis{\text{\code{log\_bs}}})$,
    90\% of the bound on those coordinates, which agrees with the pure LMetric rule on 0.924 of
    reconstructed decisions.
    % prov: Amendment A17.1
  \item \textbf{M1 with a queue threshold} (the queue-threshold ablation). M1 continued's space
    plus \code{router\_queue\_threshold} on a log scale from 0.25 to 256, starting at 32, a value
    that replays identically to the unset knob on every train cell; its baseline counterpart tunes
    the same knob on top of the best baseline's selected configuration.
    % prov: ablation (a)
\end{itemize}
% src: CR/facts/tierBC.json decisions.m2_rank2 (sources, form, warm_start), m1_continue, m1_noaff,
%      m1_v2 (bounds, sigma0 0.0028 = flip rate 0.196, inits; s3 "-e0 - 20.7 (e_log_ptok + e_log_bs),
%      90% of bound_8 (agreement 0.924 with the pure LMetric rule)"), ablation_a (q log [0.25, 256],
%      start 32, identity check), ais_league; counts "done 30"
The secondary \ais{}-informed league is trained the same way on its own build (\cref{sec:res-ais}).
% prov: A16; the league's training details are in appendix-ais-league.tex (W5).

% src: CR/facts/tierBC.json decisions.ablation_b, ablation_b_recheck_i8; CR/CONTRACT.md A17.3 (cut
%      order: ablation b first); moved from the model.tex footnote in restructure round 1 (E10)
% prov: ablation (b) = simulator-only signals
The second planned ablation, adding simulator-only signals such as engine waiting counts and KV
usage, was not run. It was first in the pre-registered order of cuts and was skipped: it needed new
plugin inputs, replay plumbing and a separate build for an arm that never ships, and free compute
existed later, so the reason was implementation cost, not compute.

\subsection{Implementing the sign constraint}

% src: CR/audits/phase2-mid/fix.md ("How (a clamp, not a box bound)", "Inits", "Equal footing");
%      CR/facts/tierBC.json decisions (m1_v2 sign_rule, m2_rank2 sign_rule, ais_league sign_rule)
\begin{itemize}
  \item \textbf{A clamp, not a box bound.} The restarts s1 and s3 start with the constrained
    coefficients of \cref{sec:sign-constraint} at 0, which on a box bound of 0 would sit on the
    edge of pycma's boundary transform, where its slope is zero and the first generation's steps on
    those coordinates shrink by about two orders of magnitude. The space therefore keeps the
    original box and clamps the decoded value, $x = \min(x, 0)$. CMA-ES searches the identical
    internal space, start, step size and seed as the unconstrained run, and generation 0 asked the
    same internal points; only the decoded policies differ.
  \item \textbf{Starts.} s1 and s3 already satisfy the constraint. The behavior clone s2 did not
    (its $\theta_3$ and $\theta_4$ were positive, so decode load attracted above about 24K tokens),
    so it was projected onto the constraint; its agreement with the cloned heuristic stayed at 0.879.
  \item \textbf{Later arms.} M1-v2 clamps every load-ordered feature (3, 4, 5, 7, 8, 10 and 12 to
    21), M2 also clamps the load columns of its context directions so that every effective load
    weight stays non-positive, and the \ais{}-informed league clamps its load and backlog features.
\end{itemize}

\subsection{Drift check before M2}
\label{sec:drift}

% src: CR/literature/LESSONS.md LR-15 (Tomlinson footnote 3; drift check); CR/facts/STATE.md
%      phase2-prep ("4 LR-15 drift runs at 96 evals"); CR/facts/gate.json full_plan.tiers
A context term can help only if the best coefficients change with the regime, and a larger model
that nests a smaller one can still lose to it on a fixed optimizer budget; Tomlinson and Benson
note that LCL, although it nests the conditional logit, can fail to beat it within a training
budget of 500 epochs \citep{tomlinson2021lcl}. The study therefore measured drift before relying
on context.
% src: CR/literature/pdfs/dcm-context/tomlinson2021-lcl-dlcl.pdf footnote 3 (pdf p.11): "LCL does not
%      beat MNL within the 500 training epochs"
% src: CR/runs/phase2/MANIFEST.json drift-{l1,l3,n4,n8} (96 evals, --val-every 3, cell/val filters);
%      CR/facts/DEVIATIONS.md 2026-10-03 "LR-15 drift runs use 96 evaluations"; p2orch.py drift_matrix
%      (val k3-10, loss = matrix[b][b] - matrix[a][b], within_mde rule, mde 0.038);
%      train subset sizes from CR/cells/train.jsonl (L1 6, L3 8, N4 18, N8 16 cells);
%      CR/facts/gate.json full_plan.cut_order_if_over[4]; CR/CONTRACT.md A17.3; LESSONS.md LR-15

\paragraph{Procedure.}
Before M2, M1 is tuned at a quarter of the budget on each of four train subsets: L1 cells only
(6 cells), L3 only (8), $\Nw = 4$ only (18) and $\Nw = 8$ only (16). The quarter budget is 96
evaluations, six whole generations, because 100 is not a multiple of 16. Each run validates every
3 generations on the matching validation cells and selects on $\krep = 0,1,2$. Each selected model
is then scored on every regime's validation cells with fresh replicates $\krep = 3,\dots,10$. The
cross-play loss of the model tuned on regime $a$, evaluated on regime $b$, is the mean clipped
log-ratio against $\piref$ of $b$'s own model on $b$'s validation cells minus that of $a$'s model;
the pairs are L1 and L3, and $\Nw = 4$ and $\Nw = 8$. Cross-play losses beyond the minimum
detectable effect argue for context, and the drifting statistic is the natural source. M2 was to
be deprioritized if all four losses lay within the MDE of 0.038. M2 at rank 2 was later made
mandatory, so the check no longer decides whether M2 trains; it is reported as evidence for or
against context terms.
% prov: "(LR-15)"; "The gate planned to deprioritize M2"; "Amendment A17 later made M2 at rank 2 mandatory"

\paragraph{Outcome.}
% src: CR/facts/tierA.json drift (cross_play_losses: l1-tuned on l3 0.0865; interpretation: verdict
%      INCONCLUSIVE, evidence_of_drift false, use_for_m2_context_sources false, why[0..3]: SE 0.0645 over
%      3 cells, distance moved 0.08-0.28 vs 0.83, pooled beats every specialist);
%      CR/facts/tierBC.json decisions.m2_rank2 (sources, why_sources)
The check came out inconclusive. The one cross-play loss above the MDE, 0.087 for the model tuned
on L1 cells evaluated on L3 cells, rests on 3 cells and is 1.3 standard errors from zero, and a
single model tuned on all cells beat every regime specialist on the specialist's own cells; the
specialists, at a quarter of the budget, had moved only a third as far from their start as the full
M1. It is therefore read as an underpowered check with no evidence of drift, and it was not used to
choose M2's sources. Those were fixed before any M2 evaluation: \code{isl\_k}, because the
documented M1 failure axis is prompt length (\cref{sec:gaming}), and
\code{mean\_active\_prefill\_tokens\_k}, the candidate set's prefill pressure.

\subsection{Compute and numerics}
\label{sec:compute}

% src: CR/CONTRACT.md A9, A10; CR/facts/phase2_prep.json targets_at_start (10 remote nodes plus the
%      workstation); CR/facts/STATE.md phase2-prep (early generation walls on EPYC 9654P and EPYC 7702P
%      nodes), tier-B/C relays i6-i9 (about 300 s and about 610-690 s per generation; ingest counts),
%      phase2-mid fix (constrained M1 re-runs 4 h 24 min to 4 h 31 min); CR/audits/phase2-mid/fix.md
Replays run on CPUs: locally through a pool of 20 slots, and on a pool of shared CPU nodes, one
tuning run per node; the first wave started on 10 remote nodes and the workstation. A generation
took 225 to \SI{333}{s} early in the main tuning, and about \SI{300}{s} later, on a 96-core EPYC
9654P node, and 610 to \SI{690}{s} on a 64-core EPYC 7702P node; the three sign-constrained M1
re-runs took 4\,h~24\,min to 4\,h~31\,min each on the latter. The main tuning ran 66 restarts of
400 evaluations in total (M1 and every tuned baseline, the later learned arms, the ablation and
the \ais{}-informed league), and each run whose ingest the campaign record logs stored between \num{30199}
and \num{30792} result records. The single test pass added \num{5220} records on the tuning build
and 900 on the \ais{} build, and the robustness pass \num{23940} (\cref{sec:res-robustness}).
Replay-seconds for the main tuning were not totaled.
% prov: "early in phase 2"; "Phase 2 ran 66 restarts ... (tiers A and B, the ablation and the AIS
%      league)" (tier A 36 incl. the three constrained M1 re-runs, tier B 12, tier C 6, AIS 12);
%      "the stage log"; "Replay-seconds for phase 2"
% src: CR/facts/test_results.json execution.evaluated_once (5,220 + 900); CR/facts/robustness.json
%      execution.records (lag 10,800 + lag-0 900 + timing 4 x 3,060 = 23,940; coverage[*].records sum to the same)

% src: CR/facts/remote.json summary, parity.totals (6,312 records, 48,901,384 rows), caveats[3]
%      (numerics host-dependent; pins equal workstation defaults; 60/60 generations),
%      parity.lr_train.numerics_check (EPYC 9654P default: first diff at generation index 0; pinned,
%      and EPYC 7702P: 60/60 identical; fake objective, budget 720, popsize 12), parity.lr_train.pinned_on_epyc_9654p
%      (best.json, best_policy.yaml, identity.json identical; history identical except gen_wall_s);
%      environment.lr_train_numerics_pin; WT/benchmarks/learned_routing/remote/node_train.sh:73-76
Runs move between machines, so a run must continue bit-exactly on any host. Two things depend on
the host: replay results depend on the operating system's math library, and CMA-ES's own linear
algebra depends on numpy's and OpenBLAS's CPU dispatch. \Cref{app:parity} gives the parity check of
every node image against the workstation (\num{6312} remote records with every field matching,
\num{48901384} per-request rows byte for byte) and the three environment variables that pin the
dispatch. With default settings, an EPYC 9654P node diverged from the workstation in the last bits
from the first generation of a 60-generation tuning run on a synthetic objective without replays;
with the pins, whose values equal the workstation's defaults, it matched in all 60 generations, and
two short real tuning runs on that node reproduced the workstation's history (apart from wall-clock
fields), selection and policy file exactly.
% prov: lr-train; best.json and best_policy.yaml are the selection and policy file.

% src: train.py module docstring ("checkpoint.pkl is written twice per generation ... stores the pycma
%      object and numpy's global RNG state"), Trainer._load_state (identity check); CR/facts/STATE.md
%      audit(build/harness-robustness-splits r0) ("lr-train kill -9 x3 resume bit-identical (history,
%      best, CMA state, continuation)")
The harness writes its checkpoint twice per generation, after sampling and after the update,
including the CMA-ES object and numpy's random state, so a resumed run continues exactly as an
uninterrupted one would; an audit killed a run three times and found its history, selection,
CMA-ES state and continuation identical (\cref{app:determinism}). A run refuses to resume if its
space, cells, reference, objective or seeds have changed.

% src: CR/facts/gate.json full_plan.cut_order_if_over[6]; CR/facts/pilot.json
%      measurements.subset_vs_full_objective (spearman 0.947, 16 candidates)
If compute ran short, the last resort in the pre-registered order of cuts was to score each
generation on a rotating 12-cell subset of the train cells, for every policy of a wave alike. On the
16 candidates of one pilot generation, subset and full-train objectives had a Spearman correlation
of 0.947 (\emph{preliminary (pilot, train subset)}).
% prov: "the gate's cut order"
